
%===================================================================================================%

\begin{frame}{The 2-$d$ Ising model}
We are going to practice data cleaning and validation in the context
of the 2-$d$ Ising model, a model for a magnet.
\begin{itemize}
\item One of the simplest magnet models
\item Analytically solvable (Cross Check 2)
\item Applications to phase transitions
\item Pedagogical for phase transitions
\end{itemize}
\end{frame}

\begin{frame}{The 2-$d$ Ising model}
\centering
\includegraphics[width=0.8\textwidth]{dataAnalysis/isingModel.pdf}
$$
H=-J\sum_{\langle ij\rangle} \sigma_i \sigma_j
~~~~~ ~~~~~ ~~~~~
\textrm{pr}(C)\sim e^{-H(C)/T}
$$
\end{frame}

\begin{frame}{Solving Ising with MCMC}
Common strategy: Generate $\nconf$ configurations randomly, then
estimate expectation value of observable $X$
$$
  \bar{X}=\frac{1}{\nconf}\sum_{i=1}^{\nconf} X_i
$$
Each $X_i$ measured on conf $C_i$ from Markov Chain
$$
  C_0\to C_1\to C_2\to...
$$
\end{frame}

% unbiased variance
% assumes uncorrelated data, but they are actually correlated
% can show variance of correlated data is larger than if one
%  had estimated it assuming independence; difference is integrated
%  autocorrelation time
% can therefore think of tau_int as factor between effectively
%  independent configurations and actual number of configurations
% can use this to correct error or to guide strategies to prune data
% note that correct estimate of tau_int demands requires preserving
%  Monte Carlo order!

% an observable you might want to estimate: magnetization
% why? It's the "order parameter" --> properties of |m|/site

% DESCRIBE TASK
% open-ended question: what kind of checks can you build in
%  while you're cleaning the data?
